\section*{第 \ref{ch:g4:series-circuits} 章——串联电路}

\begin{solution}{exo:g4:series-circuits:1}
玩家们一个接一个地坐在同一个单环上，
中间没有分叉——就像穿在同一根项链上的珠子。
\end{solution}

\begin{solution}{exo:g4:series-circuits:2}
全有或全无（一处断开，所有人都停下）；
顺序无关紧要（交换玩家，结果不变）；
分担（环上\omterm{def:g3:first-electric-circuit:bulb}{灯泡}越多，每个越暗）。
\end{solution}

\begin{solution}{exo:g4:series-circuits:3}
\omterm{def:g3:first-electric-circuit:bulb}{灯泡} 1 也会熄灭。空灯座相当于单环上的一处断口，
而第一条规则——全有或全无——让所有玩家同时停下。
\end{solution}

\begin{solution}{exo:g4:series-circuits:4}
没有任何变化：\omterm{ex:g3:first-electric-circuit:switch}{开关}仍然同时控制两个\omterm{def:g3:first-electric-circuit:bulb}{灯泡}，与之前完全一样。
第二条规则是环上的顺序无关紧要。
\end{solution}

\begin{solution}{exo:g4:series-circuits:5}
两个\omterm{def:g3:first-electric-circuit:bulb}{灯泡}时，每个都明显比只有一个时暗；
三个\omterm{def:g3:first-electric-circuit:bulb}{灯泡}时，还要更暗。\omterm{def:g3:first-electric-circuit:battery}{电池}的推力
被分担到整个环路上。
\end{solution}

\begin{solution}{exo:g4:series-circuits:6}
首尾相接：一个的 $+$ 接触下一个的 $-$，两者都沿环路朝同一方向推。
画法：$-\,[{+}\;{-}]\,[{+}\;
{-}]\,+$ 排成一行。
\end{solution}

\begin{solution}{exo:g4:series-circuits:7}
从 $+$ 出发：经过\omterm{def:g3:first-electric-circuit:battery}{电池}盒内的金属片到达按钮\omterm{ex:g3:first-electric-circuit:switch}{开关}，
经过\omterm{ex:g3:first-electric-circuit:switch}{开关}到达\omterm{def:g3:first-electric-circuit:bulb}{灯泡}，再经过\omterm{def:g3:first-electric-circuit:bulb}{灯泡}内发光的灯丝，
然后沿管子返回，经过第二节\omterm{def:g3:first-electric-circuit:battery}{电池}（首尾相接），到达第一节\omterm{def:g3:first-electric-circuit:battery}{电池}的 $-$ 极。
这是一个单环，没有分支。
\end{solution}

\begin{solution}{exo:g4:series-circuits:8}
同一环路上的任何位置都可以，例如第一个按钮之后。
\omterm{def:g4:series-circuits:series}{串联}时，两个按钮都必须闭合，才能接通这条唯一的通路，
所以只有两个按钮同时按下，铃才会响。
\end{solution}

\begin{solution}{exo:g4:series-circuits:9}
把第一个\omterm{def:g3:first-electric-circuit:bulb}{灯泡}换成已知完好的\omterm{def:g3:first-electric-circuit:bulb}{灯泡}；若灯串亮起，
说明被拿下的\omterm{def:g3:first-electric-circuit:bulb}{灯泡}就是罪魁祸首。若不亮，
就把原\omterm{def:g3:first-electric-circuit:bulb}{灯泡}装回，换到第二个灯座，依次向下检查。
在某个灯座处，灯串会亮起 ——
刚才拿下的那个\omterm{def:g3:first-electric-circuit:bulb}{灯泡}，就是坏掉的那颗珠子。
\end{solution}

\begin{solution}{exo:g4:series-circuits:10}
在黑暗中，汤姆把一节\omterm{def:g3:first-electric-circuit:battery}{电池}装反了。
头对头时，两节\omterm{def:g3:first-electric-circuit:battery}{电池}的推力相互对抗，不再接力，
\omterm{def:g3:first-electric-circuit:bulb}{灯泡}得不到任何电流——两股等力的推力相持，仍是僵局。
把其中一节\omterm{def:g3:first-electric-circuit:battery}{电池}调个头，首尾相接的接力就恢复了，灯也就亮了。
\end{solution}

\begin{solution}{exo:g4:series-circuits:11}
她说得对：在单环上，任何位置的\omterm{ex:g3:first-electric-circuit:switch}{开关}都能控制整个电路，
顺序无关紧要。有分支时，电流可以绕过她的\omterm{ex:g3:first-electric-circuit:switch}{开关}，
改走另一条支路，电动机可能仍会继续转动；
因此有分支的电路需要更仔细地安排\omterm{ex:g3:first-electric-circuit:switch}{开关}的位置，
它们是明年要学的并联电路。
\end{solution}
